\section{System overview and management model}
\label{sec:system}

HMRS is a self-hosted knowledge document management system for converting
heterogeneous documents into managed, searchable collections. A Python
parsing service, a TypeScript web application and an MCP tool server share
local storage and retrieval engines \cite{mcp}. The design follows from the
dependency chain of Section~\ref{sec:intro}: what a document is parsed as
determines what can be indexed, the collection it is filed into determines
which tier can see it, and the tiers that follow determine whether a
supporting passage is ever read. The primary object of the design is therefore
the managed collection rather than the query. State is made auditable through
three sources---the document files, the tree registry and per-library
metadata---so that a management operation can be checked against what it
claimed to do. Each check belongs to the platform operation or calling
workflow that performs it, and the evaluation reports evidence at that level.

\subsection{Management levels and the tiered access path}
\label{sec:levels}

HMRS manages documents at three levels and answers questions along a path that
mirrors them. The levels and the path are related but not identical: the
levels describe what is maintained and by which operations, while the tiers
describe what a query may look at and in what order.

\begin{table*}[t]
\centering
\small
\setlength{\tabcolsep}{4pt}
\renewcommand{\arraystretch}{1.15}
\caption{Management levels and the tiered access path. Each level names the
objects it maintains, the operations that write them, the state that describes
them, the tier that reads them, and the measurement that exercises it.
Retrieval algorithms operate only at the lower two tiers and never decide
which collection is consulted.}
\label{tab:levels}
\begin{tabular}{@{}>{\raggedright\arraybackslash}p{0.12\textwidth}>{\raggedright\arraybackslash}p{0.21\textwidth}>{\raggedright\arraybackslash}p{0.18\textwidth}>{\raggedright\arraybackslash}p{0.12\textwidth}>{\raggedright\arraybackslash}p{0.12\textwidth}>{\raggedright\arraybackslash}p{0.13\textwidth}@{}}
\toprule
Level & Objects and operations & State description & Read tier & Check & Evidence\\
\midrule
Collection & libraries and sub-libraries; automatic assignment from parsed content; description and tag derivation; catalogue listing & library registry, per-library metadata, library-level index & catalogue & catalogue inventory and description round-trip & E1, E2, E2c\\
Document & documents and their parts; duplicate review; move, rename, delete; forced re-index & document files, tree registry, per-library metadata & shelf & three-source reconciliation after each operation & E2\\
Segment & structural split; dense and lexical indices; optional evidence gate & chunk identifiers, per-segment scores with provenance & document and segment & index coverage and post-index probe & E1, E2c, E4--E8\\
\bottomrule
\end{tabular}
\end{table*}

At the collection level the system performs automatic classification: each
parsed document is matched against the catalogue of library descriptions and
assigned to the library whose scope covers its subject, with descriptions and
tags derived from the body rather than from the filename. Sub-libraries are
first-class tree nodes with their own descriptions and indices, so the
collection structure is a tree that can be consulted at query time. At the
document level the unit of work is the document and its parts; duplicate
detection reports candidate pairs for review rather than deleting them,
because adjacent parts of one source may be similar while carrying distinct
sections, and re-indexing remains a separate, explicit operation. At the
segment level a document is split at recognised structural markers and its
parts are indexed for dense and lexical search; when a workflow invokes the
evidence gate, each candidate segment is fetched and scored against the posed
question, and surviving segments are returned with provenance. That is the
only tier at which a passage can be judged to support a claim.

Figure~\ref{fig:levels} shows the two structures side by side. The
corresponding access path proceeds as a ladder: the
\emph{catalogue} tier lists collections and their descriptions, the
\emph{shelf} tier lists the documents and their descriptions inside a
selected collection, the \emph{document} tier opens a body, and the
\emph{segment} tier scores its parts. Every step narrows the scope the next
step must consider, and each step has a denominator: a catalogue that does
not contain the owning collection is a different failure from a shelf that
omits the document, and the system reports them separately.

\begin{figure*}[t]
\centering
\includegraphics[width=\textwidth]{fig_management_model}
\caption{Management levels and the tiered query path. Documents are managed
at three levels---collection, document and segment---over a shared state of
document files, the tree registry with per-library metadata, and the search
indices. A query is served by narrowing from the catalogue of collections,
through the shelf of documents in one collection, to the document body and
finally to the segment. Retrieval algorithms operate at the lower two tiers
only, and never decide which collection is consulted.}
\label{fig:levels}
\end{figure*}

\subsection{Design principles: three management contracts and enforcement layers}
\label{sec:req}
\label{sec:principles}

The three management contracts are hypotheses about where the managed
collection should govern the workflow, implemented at different layers: some
checks are mechanically enforced by server-side tools, whereas classification
and answer abstention also depend on the calling workflow, which is where the
hypotheses are weakest. We therefore distinguish the mechanism that executes
each step from the audit or evaluation that checks its outcome, and
Section~\ref{sec:eval} reports the evidence level each check reaches.

\begin{description}
\item[Content classifies the collection.] \emph{Mechanism:} an ingest agent
reads the parsed content together with the library catalogue and selects a
parent or sub-library; storage APIs validate identifiers and required fields,
not the semantic correctness of that choice. \emph{Check:} descriptions and
tags are re-read against the body, and three-source consistency is checked
after lifecycle operations (Sections~\ref{sec:acquisition}
and~\ref{sec:organization}).
\item[A tier narrows before it answers.] \emph{Mechanism:} the access path
narrows from catalogue to shelf to document to segment; when invoked, the
server-side Laya judging tool fetches or accepts document text, scores its
segments, and returns survivors with provenance; lane workflows are
instructed to treat rank as advisory. \emph{Check:} failed fetches or scores
are not admitted, and the ablation in Section~\ref{sec:e8} tests a bounded
frozen-pool contrast (Sections~\ref{sec:lanes} and~\ref{sec:gate}).
\item[Scope is reported, not assumed.] \emph{Mechanism:} the answer workflow
instructs the agent to report not-found when completed retrieval provides no
supporting evidence and to state its searched scope. This is not a
server-side semantic classifier. \emph{Check:} out-of-corpus probes and
unanswerable questions were observed as abstentions in the reported trials
(Section~\ref{sec:e5}).
\end{description}

Operationally, HMRS aims to expose parse failures, make disagreement among
its three state sources detectable, require workflow-level content inspection
before evidence use, and run on one workstation. Platform operations enforce
selected mechanical checks; agent workflows perform quality review,
classification and scoped absence reporting. Neither layer guarantees
semantic placement or answer faithfulness. The ten-domain audit probes the
collection-selection contract with a scripted route-card selection
(Section~\ref{sec:routing-audit}), where an unlogged fallback may still use
descriptions; cross-lingual membership provides a separate
filename-independent placement check (Section~\ref{sec:e1}). The verifier's
ordering role failed on the industrial pools while its filtering benefit was
not established (Sections~\ref{sec:e4}--\ref{sec:e8}).

\subsection{Architecture and subsystems}
\label{sec:architecture}

Figure~\ref{fig:arch} shows the service boundaries; Figure~\ref{fig:workflow}
traces the document and query paths. Browser requests pass through Nuxt to the
FastAPI parse service and managed MinerU worker. Agents and command-line
clients use the stdio MCP server, whose tools call web/backend APIs or read
files directly. Services share document files, the tree registry,
per-library metadata and search indices. Parse, re-index and graph operations
return pollable task identifiers. Classification, gate invocation and scoped
absence reporting remain workflow decisions; timeout, fetch and incomplete
call states remain execution failures.

\begin{figure*}[t]
\centering
\includegraphics[width=\textwidth]{fig1_architecture}
\caption{HMRS service boundaries and shared state. Browser requests pass
through the Nuxt server proxy to the FastAPI parse API and managed MinerU
worker; agents use the stdio MCP server (95 registered tools). Document
storage, tree and per-library metadata, and search indices are shared across
services. Saving and indexing are distinct steps; the Laya evidence gate is a
separate tool call and is invoked only by workflows that request it.}
\label{fig:arch}
\end{figure*}

\begin{figure*}[t]
\centering
\includegraphics[width=\textwidth]{fig2_workflow}
\caption{Document maintenance and question answering share the same
knowledge state. During ingestion, an agent reads parsed content and uses
library descriptions to choose a destination. Storage and indexing are separate
state transitions: some MCP save paths start vector/BM25 indexing
asynchronously, whereas the web parsed-file save path does not; a workflow can
then probe retrievability. At query time,
Search, Librarian and composed Hybrid workflows gather candidates and may
invoke the separate Laya gate. A scoped not-found report follows only a
completed search with no supporting evidence; timeout, fetch and tool failures
are shown as incomplete execution.}
\label{fig:workflow}
\end{figure*}

The management path covers acquisition (Section~\ref{sec:acquisition}),
content-informed organisation and lifecycle checks
(Section~\ref{sec:organization}), and the tool and runtime layer
(Section~\ref{sec:implementation}). The chat endpoint defines pinned-search
and all-library librarian policies. Hybrid is a composed retrieval
configuration, not a third built-in server policy; it combines vector and
catalogue calls under an explicit directive. The document graph supports
navigation, not evidence. Experience retrieval and the persona-distillation workflows reuse
knowledge-base infrastructure but follow their own retrieval and filtering
paths; neither is described here as invoking Laya, and neither was evaluated
at scale.

The platform exposes five management domains: acquisition, organisation,
lifecycle, query and operational experience. Some controls are enforced by
API/MCP operations; semantic and quality checks are agent-workflow procedures,
and campaign audits test selected outcomes. Table~\ref{tab:functions}
distinguishes these layers and maps each domain to the evidence that exercises
it.

\begin{table*}[t]
\centering
\small
\setlength{\tabcolsep}{4pt}
\caption{Management functions, the enforcement or workflow check associated
with each, and the evaluation that exercises it. The table separates API/MCP
controls from agent procedures and campaign audits; semantic choices such as
which collection owns a document remain workflow decisions and are measured,
not enforced.}
\label{tab:functions}
\begin{tabular}{@{}>{\raggedright\arraybackslash}p{0.13\textwidth}>{\raggedright\arraybackslash}p{0.37\textwidth}>{\raggedright\arraybackslash}p{0.30\textwidth}>{\raggedright\arraybackslash}p{0.11\textwidth}@{}}
\toprule
Group & Operations & Platform or workflow check & Evidence\\
\midrule
Acquisition & submit, parse (MinerU), quality review, duplicate review, structural split, store, index & parse-task state; workflow quality checks; entry-point-specific deduplication and indexing; post-index probe & E1\\
Organisation & create library or sub-library, content-informed placement, descriptions and tags, duplicate review & API validates identifiers and required fields; workflow checks descriptions and tags; audit compares files, tree and library metadata & E1--E2\\
Lifecycle & move, rename, delete, force re-index, graph build & operation-specific API updates; reconciliation audit detects selected stale or inconsistent state & E2\\
Query & tier selection, vector, keyword, catalogue and two-stage search, gate invocation & endpoint enforces tool allow/deny lists, timeouts and turn limits; invoked gate returns survivor provenance & E4--E8\\
Experience & extract, approve, tier, search & experience tools apply their own retrieval, verification and credibility filters & E3\\
\bottomrule
\end{tabular}
\end{table*}

\subsection{Data model and object lifecycle}
\label{sec:datamodel}
\label{sec:lifecycle}

Documents are Markdown with extracted images. Library state is described by
the document files, the \texttt{.tree-fs.json} registry and each
\texttt{.knowledge-base.yml}. The ingest workflow tracks raw, parsed, stored
and indexed states and can mark completion after its configured retrieval
probe. Indexing is a separate transition: some save paths start vector/BM25
indexing asynchronously, while the web parsed-file save path does not.
Reconciliation detects selected stale-index conditions or disagreement among
the three sources; it is not a cross-store transaction. Because the three
sources are written by different management operations, the collection has no
single authoritative record of its own state, and the reconciliation audit
exists to make disagreement observable rather than to prevent it.
